\documentclass[11pt]{article}

\usepackage[margin=1in]{geometry}
\usepackage{amsmath,amssymb,amsthm,mathtools}
\usepackage[hidelinks]{hyperref}
\usepackage[numbers,sort&compress]{natbib}

\newtheorem{theorem}{Theorem}
\newtheorem{lemma}{Lemma}
\newtheorem{corollary}{Corollary}
\newtheorem{proposition}{Proposition}
\newtheorem{definition}{Definition}
\newtheorem{remark}{Remark}

\newcommand{\cN}{\mathcal N}
\newcommand{\cC}{\mathcal C}
\newcommand{\cH}{\mathcal H}
\newcommand{\cA}{\mathcal A}
\newcommand{\cD}{\mathcal D}
\newcommand{\cE}{\mathcal E}
\newcommand{\cM}{\mathcal M}
\newcommand{\cR}{\mathcal R}
\newcommand{\cT}{\mathcal T}
\newcommand{\cZ}{\mathcal Z}
\newcommand{\Tr}{\operatorname{Tr}}

\title{Horizon Interfaces and the Recovery of Private Quantum Information}
\author{}
\date{\today}

\begin{document}

\maketitle

\begin{abstract}
Horizons behave like quantum/classical interfaces: coarse no-hair data
are public and classical, while microscopic or diary information remains
private until the observer's access algebra is enlarged.  This note
formalizes that interface in finite-dimensional quantum-information
terms.  In exact form, redundant recovery from disjoint record fragments
forces the public algebra to be commutative, giving an effective center
and private noncommuting blocks; approximate public records are treated
as the standard Darwinism/SBS stability input.  In any finite-velocity
source-local realization, a diary deposited far from the access region is
absent from the reduced emitted record, up to the Lieb-Robinson tail,
before the routing time; no decoder on those records can recover it in
that window, even with unbounded computational power.  This is an
information-theoretic obstruction, independent of decoding complexity.
Source-side
saturation does not remove the latency, and anonymous unitary saturated
couplings can record only a collective center while leaving private
blocks inaccessible.  The exact protected private complement is the
commutant/noiseless subsystem of the record-generating algebra.
The resulting recovery problem has two independent bottlenecks: access
geometry and coherent export.  A horizon-like interface with redundant
public records, private blocks, saturated source access, and
location-uniform fast private recovery must use fast internal
routing/scrambling or nonlocal/dressed access; conversely, a
log-diameter access geometry equipped with a theorem-backed
decoupling/export mixer realizes the fast-routing branch in
logarithmic or polylogarithmic depth.  A second-moment export criterion
then reduces the deterministic sparse-dynamics strengthening to a
specific moment-gap problem.  Freezing the internal dynamics separates
routed recovery from dressed access operationally.
\end{abstract}

\section{Introduction}

The motivating distinction is between public classical data and private
quantum data.  Public data may be redundantly available in many records:
pointer positions, thermodynamic labels, or no-hair charges.  Private
data are noncommuting block degrees of freedom inside a fixed public
sector, such as a diary state or a black-hole microstate.  Quantum
Darwinism and no-broadcasting explain why redundant objective records
select commuting public information \citep{Barnum1996NoBroadcasting,
Brandao2015GenericClassical,Le2019StrongQD,Korbicz2020Roads}.  The
same access-relative distinction is standard in quantum cryptography:
secret-sharing schemes distinguish authorized from unauthorized record
sets, data-hiding schemes make information globally present but
inaccessible to restricted observers, private quantum channels randomize
states without a key, and locking shows how side information can sharply
change accessible information
\citep{CleveGottesmanLo1999,TerhalDiVincenzoLeung2001,
EggelingWerner2002,AmbainisMoscaTappDeWolf2000,
DiVincenzoHaydenTerhal2004}.  The question addressed here is not the
existence of such access structures in the abstract.  It is when a
physically specified record algebra, generated by local dynamics,
permits private quantum information to become recoverable quickly from
the allowed records.

The answer is a mechanism classification.  For finite-velocity
source-local systems, private information cannot be recovered before it
has routed to the source-local access region.  Thus horizon-like fast
private recovery requires either fast internal routing/scrambling, as in
Hayden-Preskill-type pictures \citep{Hayden2007BlackHoles}, or a
nonlocal/dressed access algebra, as in gravitationally constrained or
holographic descriptions.  The obstruction is information-theoretic:
allowing an arbitrarily powerful decoder on the emitted records does not
help before the diary has reached the record channel, unless the access
algebra already contains nonlocal or dressed diary information.

The positive side is equally important.  Once the diary reaches an
export-active block, recovery requires coherent export: the emitted
records must decouple the diary reference from the remaining source.
This yields an access/export window.  Ordinary local reservoirs have a
power-law access time for worst-case deposits, while bounded-degree
log-diameter access geometries can have logarithmic access time.  With a
theorem-backed design, tensor-product-expander, or random-circuit
decoupling primitive, such a geometry realizes the fast-routing branch.
For a deterministic sparse dynamics, the corresponding target is a
second-moment gap in the decoupling functional for the emitted-record
partition.

The claim is about finite-dimensional access models and their chosen
factorizations.  It does not derive gravitational dynamics.  The
gravitational language enters when a horizon is modeled as an interface
with public no-hair data, private microscopic blocks, and an enlarged or
dressed observer algebra.

The public side is deliberately not the novel part of the paper.  The
exact cut theorem below records the no-broadcasting core, while
approximate public objectivity is imported from Quantum Darwinism and
spectrum-broadcast-structure assumptions.  The new constraints concern
the private complement: when it is absent from accessible records, when
it can become recoverable, and what mechanisms can make that recovery
fast.  The algebraic obstruction is the commutant of the
record-generating interaction algebra; the dynamical question is how
quickly that commutant becomes trivial inside a public block and whether
the accessible records actually export the diary rather than merely
becoming sensitive to it.

\section{The Cut Theorem}

The public/private split has an exact finite-dimensional form.  Let
$\cA\subset B(\cH)$ be the algebra of information to be made available
in records.  Say that $\cA$ is redundantly recoverable from two disjoint
record fragments $R_1$ and $R_2$ if the record-formation channel
$\cE:\mathsf S(\cA)\to \mathsf S(R_1R_2F)$ has reduced channels
$\cE_i=\Tr_{\overline{R_i}}\circ\cE$ and there are recovery channels
$\cD_i:\mathsf S(R_i)\to\mathsf S(\cA)$ such that
\begin{equation}
  \cD_i\circ \cE_i = {\rm id}_{\mathsf S(\cA)},
  \qquad i=1,2 .
\end{equation}

\begin{theorem}[Redundant recoverability is central]
\label{thm:cut}
If a finite-dimensional algebra $\cA$ is exactly redundantly recoverable
from two disjoint record fragments, then $\cA$ is commutative.
\end{theorem}

\begin{proof}
Every finite-dimensional $C^*$-algebra has the form
\begin{equation}
  \cA \cong \bigoplus_x B(\cH_x).
\end{equation}
If some block has $\dim \cH_x>1$, then redundant recovery from $R_1$ and
$R_2$ would broadcast the full set of quantum states on $B(\cH_x)$.  The
channel $(\cD_1\otimes\cD_2)\circ\Tr_F\circ\cE$, restricted to that
block, has both marginals equal to the input state.  This contradicts
the no-broadcasting theorem \citep{Barnum1996NoBroadcasting}, since the
states in a nontrivial matrix block do not commute.  Hence every block
is one-dimensional, so $\cA$ is commutative.
\end{proof}

\begin{corollary}[Public center and private blocks]
\label{cor:public-private}
Let the logical algebra have block form
\begin{equation}
  \cA_{\rm code}
  =
  \bigoplus_x B(\cH_x),
  \qquad
  \cZ(\cA_{\rm code})
  =
  \operatorname{span}\{P_x\}.
\end{equation}
The sector label $x$ may be redundantly public.  The noncommuting
degrees of freedom inside each block cannot be redundantly recoverable
from disjoint fragments.  They are private to those fragments unless an
enlarged or differently dressed access algebra is supplied.
\end{corollary}

This corollary is the public/private algebraic version of a standard
access-structure statement: small or disjoint fragments may be
unauthorized for noncommuting information even when they redundantly
carry a classical label.  The role of the later sections is to add the
physical constraints absent from abstract secret-sharing or data-hiding
definitions: locality of the record coupling, routing latency,
source-side participation, coherent export, and the possibility of
dressed access.

In the exact block-scalar case this split is visible directly.  If a
fragment channel has the form
\begin{equation}
  \cN_a(\rho)
  =
  \sum_x \Tr(P_x\rho)\,\sigma_x^{(a)},
\end{equation}
then every fragment observable $B_a$ pulls back to
\begin{equation}
  \cN_a^*(B_a)
  =
  \sum_x \Tr(B_a\sigma_x^{(a)})\,P_x
  \in \cZ(\cA_{\rm code}).
\end{equation}
The fragment sees only the central distribution $p_x=\Tr(P_x\rho)$.
All states with the same $p_x$, including all diary states inside a
fixed block, are identical to that fragment.

The paper uses this exact result as the algebraic core of the public
side.  Approximate Darwinian or spectrum-broadcast records require an
additional stability input: many fragments identify $x$ with small
error, while within-block distinguishability remains small below the
relevant access threshold.  We import that approximate objectivity
condition from the Quantum Darwinism and SBS literature
\citep{Brandao2015GenericClassical,Le2019StrongQD,Korbicz2020Roads};
the quantitative latency results below do not depend on a particular
choice of stability constants.

\section{Source-Local Record Model}

Let $\Lambda$ be a bounded-degree graph or lattice with bounded local
Hilbert dimension and finite-range, uniformly bounded Hamiltonian
\begin{equation}
  H_\Lambda=\sum_Z h_Z,
  \qquad
  \sup_x \sum_{Z\ni x}\lVert h_Z\rVert < \infty .
\end{equation}
Let $X\subset\Lambda$ be the access region and $Y\subset\Lambda$ the
region where a diary is deposited, with $L=d(X,Y)$.  The record system
is an environment or field coupled to the source only near $X$:
\begin{equation}
  H_{\rm tot}(s)=H_\Lambda+H_R(s)+H_{\rm int}(s),
  \qquad
  \operatorname{supp}_\Lambda H_{\rm int}(s)\subset N_R(X).
\end{equation}
The emitted record $R_t$ is the part of the environment collected up to
time or record depth $t$.  The diary channel is
\begin{equation}
  \cN_t(\rho_D)
  =
  \Tr_{\neg R_t}\!\left[
  W_t(\rho_D\otimes\rho_{\rm rest}\otimes |0\rangle\langle 0|_R)
  W_t^\dagger
  \right],
\end{equation}
where $W_t$ is generated by $H_{\rm tot}(s)$.

\begin{definition}[Source-local emission]
An emitted record channel is source-local relative to $X$ if the source
support of the interaction that creates the record remains in a bounded
neighborhood of $X$ throughout the record window.
\end{definition}

\begin{definition}[Location-uniform fast recovery]
A system has location-uniform fast recovery for $k$-qubit deposits if
diaries deposited in arbitrary regions $Y\subset\Lambda$ are recoverable
from the allowed record/access algebra in sub-power-law record depth or
time, with the stated side-information and decoder resources.
\end{definition}

\section{Reduced-Record Locality Bound}

\begin{lemma}[Reduced-record Lieb-Robinson bound]
\label{lem:reduced-lr}
In the source-local record model, let $\cC_t$ be the constant channel
that maps every diary state to the record state produced by a fixed
reference diary preparation.  For $t<L/v$,
\begin{equation}
  \lVert \cN_t-\cC_t\rVert_\diamond
  \le
  \varepsilon_{\rm LR}(t),
  \qquad
  \varepsilon_{\rm LR}(t)
  \le
  C |X|\,|Y|\, t\, e^{-\mu(L-vt)} ,
\end{equation}
with constants determined by local dimensions, interaction range,
operator-norm bounds, and the chosen Lieb-Robinson estimate.  For a
$k$-qubit diary supported on bounded-degree cells, $|Y|=O(k)$.
\end{lemma}

\begin{proof}[Proof note]
Discretize the record field as environment ancillas coupled only near
$X$, if needed.  A perturbation or diary operation supported in $Y$ can
influence the reduced state of the collected record only through the
combined source-plus-environment path $Y\to X\to R_t$.  Since
$H_{\rm int}$ has source support only near $X$, there is no direct
$Y\to R_t$ shortcut.  The standard information-propagation consequence
of Lieb-Robinson bounds \citep{Lieb1972FiniteVelocity,Bravyi2006LR}
gives the displayed exponential tail after summing over the spacetime
boundary of the coupling region and over the diary support.

For the diamond norm, include an arbitrary diary reference as a
spectator.  Equivalently, test the channel with controlled diary
operations
\begin{equation}
  W_{YA}=\sum_j |j\rangle\!\langle j|_A\otimes U_Y^{(j)}
\end{equation}
on the diary support and an ancilla $A$.  Record observables commute
with $A$, and the Lieb-Robinson estimate applies uniformly to each
$Y$-supported block $U_Y^{(j)}$.  Taking the supremum over controlled
unitaries gives the same reduced-record bound for entangled
diary-reference inputs, hence the displayed diamond-norm estimate.
\end{proof}

\begin{lemma}[Constant channels do not recover diaries]
\label{lem:constant}
If $\lVert \cN_t-\cC_t\rVert_\diamond\le \varepsilon$, then for every
decoder $\cR:R_t\to \widehat D$,
\begin{equation}
  F_e\!\left[
  (\operatorname{id}_{R_D}\otimes \cR\cN_t)(\Phi_{R_DD}),
  \Phi_{R_D\widehat D}
  \right]
  \le
  \frac{1}{d_D}+O(\varepsilon).
\end{equation}
\end{lemma}

\begin{proof}
Post-processing by $\cR$ cannot increase diamond distance.  The constant
channel produces a product state across $R_D|\widehat D$, whose overlap
with a $d_D$-dimensional maximally entangled target is at most $1/d_D$.
Entanglement fidelity is a linear overlap with the target state, so
trace-norm error changes it by $O(\varepsilon)$.
\end{proof}

\section{Access-Latency Classification}

\begin{theorem}[Source-local records forbid horizon-like latency]
\label{thm:latency}
For a finite-dimensional source with finite Lieb-Robinson velocity and
source-local emission, a diary deposited distance $L$ from the access
region cannot be recovered with high entanglement fidelity from the
emitted record before
\begin{equation}
  t \ge L/v - O(\log(C|X|\,|Y|t/\varepsilon)).
\end{equation}
In a $d$-dimensional local reservoir with entropy $S=\Theta(|\Lambda|)$
and worst-case deposits at distance $L=\Theta(S^{1/d})$, the
location-uniform recovery latency is power-law:
\begin{equation}
  t_{\rm rec}=\Omega(S^{1/d})
\end{equation}
up to logarithmic and geometry-dependent factors.
\end{theorem}

\begin{proof}
Combine Lemmas~\ref{lem:reduced-lr} and~\ref{lem:constant}.  The
logarithm contains $t$ and the diary size only through slowly varying
prefactors, so the implicit fixed point shifts the ballistic scale only
logarithmically for $k=O(\operatorname{polylog}S)$ deposits.
\end{proof}

\begin{corollary}[Mechanism classification]
Location-uniform sub-power-law recovery of arbitrary private deposits
requires at least one of:
\begin{enumerate}
  \item fast internal routing/scrambling, so the relevant dynamics has no
  finite-velocity bottleneck; or
  \item nonlocal/dressed access, so the allowed record algebra is not
  generated by source-local emission in factorized source variables.
\end{enumerate}
\end{corollary}

\begin{corollary}[Frozen-dynamics diagnostic]
\label{cor:frozen}
Consider the same source-local record model, but freeze the internal
source dynamics that transports operators from the diary region $Y$ to
the access region $X$, while keeping the record coupling near $X$.  If a
diary deposited in $Y$ remains recoverable from the allowed records with
order-one entanglement fidelity, then the access algebra used by the
decoder is not generated by source-local emission in the factorized
source variables, or the side information already contains the diary.
\end{corollary}

\begin{proof}
With the internal source dynamics frozen, source-local record formation
near $X$ has no path by which a perturbation in the disjoint region $Y$
can affect the emitted record.  Equivalently, Lemma~\ref{lem:reduced-lr}
has zero routing velocity for the $Y\to X$ path, so the diary-to-record
channel is constant on the diary, up to any explicitly supplied
nonlocal/dressed access or trivial diary-containing side information.
Lemma~\ref{lem:constant} then excludes order-one recovery from the
source-local record alone.
\end{proof}

\begin{remark}
This is an operational separator between the two fast branches.  If
fast recovery disappears when internal routing is frozen, the recovery
was routing/scrambling-assisted.  If it survives, the recovering algebra
already had nonlocal or dressed support relative to the factorized
source variables.
\end{remark}

\begin{theorem}[Source-side saturation does not imply fast recovery]
\label{thm:saturation-slow}
For each $S$ there is a finite-dimensional record channel with
source-side participation
\begin{equation}
  N_{\rm eff}(W)
  =
  \frac{(\Tr W)^2}{\Tr W^2}
  =
  S
\end{equation}
but worst-case private recovery latency $\Omega(S)$ records.
\end{theorem}

\begin{proof}
Let the source consist of $S$ cells, each carrying one private qubit.
For each cell $i$ let $J_i$ be a normalized local source operator, and
take the source-side Gram matrix
\begin{equation}
  W_{ij}=\langle J_i,J_j\rangle=\delta_{ij}.
\end{equation}
Then $N_{\rm eff}(W)=S$, so the source-side access family is saturated.

Now define the emitted record process to have only one elementary output
per step.  At each step the channel samples one cell $i$ uniformly and
applies a local record map to that cell.  The record may reveal the
sampled address or hide it; the lower bound below does not use address
hiding.  Consider a diary deposited in a worst-case unknown cell $j$.
After $n$ records, the probability that cell $j$ has ever been sampled is
at most $n/S$.

Conditioned on the event that $j$ has not been sampled, the emitted
record is independent of the diary.  Thus the $n$-record diary-to-record
channel has the form
\begin{equation}
  \cN_n
  =
  \left(1-p_n\right)\cC_n+p_n\cM_n,
  \qquad
  p_n\le \frac{n}{S},
\end{equation}
where $\cC_n$ is a constant channel on the diary and $\cM_n$ is some
channel produced on the sampled branch.  Hence
\begin{equation}
  \lVert \cN_n-\cC_n\rVert_\diamond
  \le
  2p_n
  \le
  \frac{2n}{S}.
\end{equation}
By Lemma~\ref{lem:constant}, every decoder has entanglement fidelity at
most $1/d_D+O(n/S)$.  An order-one recovery advantage therefore requires
$n=\Omega(S)$ records.
\end{proof}

\begin{remark}
The construction isolates the issue.  Source-side saturation says that
the family of source couplings has full participation.  It does not say
that a newly deposited diary reaches the emitted record quickly.  The
failure here is bandwidth/sampling-limited rather than LR-limited; the
finite-velocity slow router of Theorem~\ref{thm:latency} is the
geometric version of the same separation.
\end{remark}

\begin{remark}
The theorem is stated without useful side information about the diary
location.  Side information about unsampled cells does not alter the
bound, since those cells have not interacted with the emitted record.
Side information that already contains the diary or its purifier is the
trivial excluded case.  A Hayden-Preskill escape uses a different model:
fast scrambling first routes the diary into global degrees of freedom,
and previously collected radiation supplies the side information needed
for recovery from a small number of fresh records.
\end{remark}

\begin{theorem}[Anonymous saturation does not force mixing]
\label{thm:anonymous-saturation}
Source-side saturation, record-label anonymity, source permutation
symmetry, and unitary record formation do not imply private-block
mixing or private recovery.
\end{theorem}

\begin{proof}
Let the source be $S$ qubits and let the microscopic source operators be
the normalized local observables
\begin{equation}
  J_i=\sigma_i^z,\qquad i=1,\ldots,S .
\end{equation}
With the Hilbert-Schmidt inner product normalized on the source,
$W_{ij}=\langle J_i,J_j\rangle=\delta_{ij}$, so
$N_{\rm eff}(W)=S$.  This is source-side saturation of the microscopic
contribution family before compression into the single collective
pointer channel.

Let $P_q$ be the projector onto the eigenspace of the permutation-
invariant charge
\begin{equation}
  Q=\sum_{i=1}^S \sigma_i^z ,
\end{equation}
and let the record be a finite pointer with at least $2S+1$ orthogonal
positions and shift operator $V$.  The controlled-shift interaction
\begin{equation}
  U=\sum_q P_q\otimes V^q
\end{equation}
is unitary and permutation invariant.  Repeating the interaction with
fresh pointers gives arbitrarily many record copies of the collective
charge $Q$.  The record label contains no microscopic source address,
and the channel is covariant under permutations of the source qubits.

The record state depends only on the eigenspace of $Q$.  For a charge
eigenvalue $q$, the corresponding source subspace has dimension
\begin{equation}
  \binom{S}{(S+q)/2}
\end{equation}
whenever this expression is defined.  All density matrices supported
inside the same $Q$-eigenspace produce the same record statistics under
every repetition of this coupling.  Thus the public algebra generated by
the records is the commutative algebra of projectors onto $Q$ sectors,
while the noncommuting algebra inside a degenerate sector remains
private forever.  No mixing of private block information into the record
is forced by saturation, anonymity, symmetry, or unitarity.
\end{proof}

\section{Protected Private Complements}

The collective-charge example is not an accident.  It is the standard
noiseless-subsystem obstruction in the language of access channels
\citep{Zanardi1997Noiseless,Knill2000GeneralNoise,Kribs2005Unified}.

\begin{definition}[Record-generating algebra]
For a family of passive record instruments on a finite source Hilbert
space $\cH$, let $\cA_{\rm rec}$ be the unital $*$-algebra generated by
the Heisenberg pullbacks of all bounded record observables in the
instrument family.  Equivalently, in a microscopic interaction model it
is the algebra generated by the source operators through which the
record couples, together with their products and adjoints.
\end{definition}

\begin{remark}
The algebra is relative to the specified access family.  Enlarging the
allowed instruments, adding active probes, or including dressed
observables enlarges $\cA_{\rm rec}$ and can shrink the protected
commutant.  The theorem below therefore identifies protected private
information relative to a chosen record/access model, rather than an
absolute property of the source Hilbert space.
\end{remark}

\begin{theorem}[Protected complement is the commutant]
\label{thm:commutant}
Let $\cA_{\rm rec}\subset B(\cH)$ be the record-generating algebra.  In
finite dimension there is a decomposition
\begin{equation}
  \cH
  =
  \bigoplus_\alpha
  \cH_{\alpha}^{L}\otimes \cH_{\alpha}^{R},
  \qquad
  \cA_{\rm rec}
  =
  \bigoplus_\alpha
  B(\cH_{\alpha}^{L})\otimes I_{\alpha}^{R},
\end{equation}
and
\begin{equation}
  \cA_{\rm rec}'
  =
  \bigoplus_\alpha
  I_{\alpha}^{L}\otimes B(\cH_{\alpha}^{R}) .
\end{equation}
Record statistics generated by $\cA_{\rm rec}$ are insensitive to the
full quantum state on each factor $\cH_{\alpha}^{R}$.  Thus the exact
private complement protected forever from the passive records is
precisely the noncommuting part of the commutant $\cA_{\rm rec}'$.
\end{theorem}

\begin{proof}
The displayed decomposition is the finite-dimensional structure theorem
for unital $*$-algebras.  Every Heisenberg record effect $E$ lies in
$\cA_{\rm rec}$, so on the $\alpha$ block it has the form
$E_\alpha^L\otimes I_\alpha^R$.  For any two states with the same
left-factor reduced states
\begin{equation}
  \Tr_{\cH_\alpha^R}(P_\alpha\rho P_\alpha)
  =
  \Tr_{\cH_\alpha^R}(P_\alpha\sigma P_\alpha)
  \quad \text{for all }\alpha ,
\end{equation}
one has $\Tr(E\rho)=\Tr(E\sigma)$ for every record effect and therefore
for every measurement or decoder using only the passive record algebra.
Conversely, operators in the commutant act on the multiplicity factors
$\cH_\alpha^R$ and are not distinguished by $\cA_{\rm rec}$.  These are
exactly the noiseless/private subsystems relative to the record
instruments.
\end{proof}

\begin{corollary}[Irreducibility criterion]
\label{cor:irreducible}
Within a public block, there is no exact protected private subsystem
relative to the passive record family iff the record-generating algebra
acts irreducibly on that block, equivalently iff its commutant inside the
block is only the scalar algebra.
\end{corollary}

\begin{proof}
This is Theorem~\ref{thm:commutant} restricted to one public block.  A
nontrivial multiplicity factor $\cH_\alpha^R$ gives a protected private
subsystem; if all multiplicity factors are one-dimensional, the
commutant contributes only block scalars.
\end{proof}

\begin{remark}
The collective-charge pointer realizes the degenerate case.  The
record-generating algebra is the abelian algebra of functions of
$Q=\sum_i\sigma_i^z$, so its commutant contains the full matrix algebra
inside each fixed-charge sector.  Saturation, anonymity, symmetry, and
unitarity leave that commutant intact.  A genuine mixing condition must
therefore say that the time-generated record algebra becomes irreducible
inside the relevant public block, or quantify how rapidly its commutant
shrinks in an approximate sense.
\end{remark}

\begin{definition}[Heisenberg observability algebra]
For a fixed source Hamiltonian $H$ and a record coupling through source
operator(s) $K_a$, define the observability algebra up to sampling times
$T=\{t_1,\ldots,t_n\}$ by
\begin{equation}
  \cA_T
  =
  \operatorname{alg}\{
    e^{iHt_j}K_a e^{-iHt_j},
    e^{iHt_j}K_a^\dagger e^{-iHt_j}
    : a,j
  \}.
\end{equation}
The exact private algebra invisible to all passive records generated at
these times is $\cA_T'$.
\end{definition}

\begin{corollary}[Exact observability criterion]
\label{cor:observability}
Within a public block, no exact protected private subsystem survives the
sampled Heisenberg orbit iff
\begin{equation}
  \cA_T'=\mathbb C I,
\end{equation}
equivalently iff $\cA_T$ acts irreducibly on the block.
\end{corollary}

\begin{proof}
Apply Corollary~\ref{cor:irreducible} to the record-generating algebra
$\cA_T$.  The equivalence between irreducibility and scalar commutant is
Burnside's theorem for finite-dimensional complex matrix algebras.
\end{proof}

\begin{proposition}[Generic exact observability]
\label{prop:generic-observability}
Fix a $D$-dimensional public block and sampling times
$T=\{t_1,\ldots,t_n\}$.  For sampled operators
$L_j=e^{iHt_j}Ke^{-iHt_j}$, the condition
\begin{equation}
  \operatorname{alg}\{L_j,L_j^\dagger:j=1,\ldots,n\}
  = B(\mathbb C^D)
\end{equation}
is Zariski open in the matrix entries of the $L_j$ whenever it is
nonempty.  Consequently, in any finite-parameter family $(H,K)$ for
which one point gives a full sampled algebra and for which this
nonvanishing determinant is not identically zero, the set of parameters
with no exact protected commutant is open and dense away from the
determinant-zero locus.
\end{proposition}

\begin{proof}
A finite set of matrices generates the full matrix algebra iff words in
those matrices of length at most $D^2-1$ span the $D^2$-dimensional
vector space $B(\mathbb C^D)$.  This is equivalent to at least one
$D^2\times D^2$ determinant, formed from vectorized words, being
nonzero.  The determinant is a polynomial in the matrix entries of the
sampled operators $L_j$ and $L_j^\dagger$.  Therefore the full-algebra
condition is the nonvanishing of at least one such determinant: a
Zariski-open condition in the sampled operator variables.  Pulling this
determinant back along any chosen parameterization $(H,K)\mapsto L_j$
gives an analytic function of those parameters.  If that function is not
identically zero, its nonzero set is open and dense, and on that set the
commutant is scalar by Corollary~\ref{cor:observability}.
\end{proof}

\begin{remark}
The proposition is an exact genericity statement, not a scrambling-rate
theorem.  It says that exact protected commutants are nongeneric once
the sampled orbit is allowed to explore a full matrix algebra.  It does
not say how many samples are needed, how stable the conclusion is
approximately, or whether deterministic chaotic dynamics produces the
iid-like independence used in Theorem~\ref{thm:pauli-growth}.
\end{remark}

\begin{corollary}[Exact obstruction classes]
\label{cor:obstructions}
The sampled orbit has a nontrivial protected commutant whenever there is
a non-scalar operator $Q$ on the public block such that
\begin{equation}
  [Q,e^{iHt_j}K_a e^{-iHt_j}]=
  [Q,e^{iHt_j}K_a^\dagger e^{-iHt_j}]=0
  \qquad \text{for all }a,j .
\end{equation}
In particular, exact protection survives in any sector left invariant by
symmetries, selection rules, fragmentation, collective monitoring, or an
operator Krylov subspace that fails to generate the full matrix algebra.
\end{corollary}

\begin{proof}
Such a $Q$ lies in $\cA_T'$ and is not scalar, so
Corollary~\ref{cor:observability} gives a protected private algebra.
The listed mechanisms are standard ways for all sampled record
operators to preserve a common nontrivial commutant or invariant block
decomposition.
\end{proof}

\begin{definition}[One-step observability channel]
\label{def:obs-channel}
For a sampled record operator family $\mathcal L_i=\{L_{i,a}\}_a$ on a
public block, let $\cA_i=\operatorname{alg}\{L_{i,a},L_{i,a}^\dagger:a\}$
and let $\Pi_i$ be the Hilbert-Schmidt orthogonal conditional
expectation onto the commutant $\cA_i'$.  The one-step observability
channel on operators is
\begin{equation}
  \cT_i(O)=\Pi_i(O).
\end{equation}
For records accumulated through $n$ steps, define
\begin{equation}
  \cT_{1:n}=\cT_n\cdots \cT_1 .
\end{equation}
This channel keeps the component of an operator that remains invisible
to each sampled record family separately.  Since
\begin{equation}
  \cA_{1:n}'=
  \left(\operatorname{alg}(\cA_1,\ldots,\cA_n)\right)'
  =
  \bigcap_{i=1}^n \cA_i',
\end{equation}
the conditional expectation $\mathbb E_n$ onto the exact protected
algebra $\cA_{1:n}'$ is the projection onto the common fixed subspace of
the $\Pi_i$.
\end{definition}

\begin{definition}[Observability gap on a diary algebra]
Let $\cD_0$ be the traceless operator space of a fixed $k$-qubit diary
inside a public block.  The sampled record sequence has observability
gap certificate $\gamma>0$ on $\cD_0$ with constant $C$ if for every
$O\in\cD_0$,
\begin{equation}
  \lVert \cT_{1:n}(O)\rVert_2
  \le
  C e^{-\gamma n}\lVert O\rVert_2 .
\end{equation}
\end{definition}

\begin{remark}
The product $\cT_{1:n}$ is a sufficient certificate for approximate
de-protection, not a necessary definition of it.  For a deterministic
sequence of noncommuting one-step commutants, its contraction is an
alternating-projections problem governed by the relative angles among
the invisible subspaces.  A weak product bound therefore does not by
itself show that private information remains protected.  The physical
target is a conditional, or windowed, contraction statement for the
actual Heisenberg record sequence; the product certificate is the
cleanest way to turn such a contraction into a bound on the protected
component.
\end{remark}

\begin{lemma}[Observability gap implies de-protection]
\label{lem:obs-gap}
If the sampled record sequence has observability gap $\gamma$ on
$\cD_0$, then the exactly protected component of every diary operator
$O\in\cD_0$ obeys
\begin{equation}
  \lVert \mathbb E_n(O)\rVert_2
  \le
  C e^{-\gamma n}\lVert O\rVert_2 .
\end{equation}
Thus the fixed-diary effective irreducibility rate is at least
$\gamma$.
\end{lemma}

\begin{proof}
The exact protected algebra $\cA_{1:n}'$ is contained in every one-step
commutant $\cA_i'$.  Therefore the projection $\mathbb E_n(O)$ is fixed
by each $\Pi_i$ and lies in the common fixed space of the product
$\cT_{1:n}$.  Since $\mathbb E_n$ is an orthogonal projection onto a
subspace of that common invisible space,
\begin{equation}
  \lVert \mathbb E_n(O)\rVert_2
  \le
  \lVert \cT_{1:n}(O)\rVert_2 .
\end{equation}
The claimed bound follows from the observability-gap hypothesis.
\end{proof}

\begin{remark}
The definition is deliberately channel-level rather than exact-algebraic.
Exact observability asks whether $\cA_{1:n}'$ is scalar.  The gap asks
how quickly a fixed diary loses its component in the invisible directions
resolved by the record sequence.  Small couplings, hydrodynamic modes,
nearly conserved quantities, and correlated Heisenberg iterates can all
make $\gamma$ small even when the exact commutant is already scalar.
\end{remark}

\begin{proposition}[Benchmark observability gaps]
\label{prop:gap-benchmarks}
The observability-gap language has the following benchmark behavior.
\begin{enumerate}
  \item Suppose there is an adapted sequence of sampled algebras and a
  number $\eta>0$ such that, at step $i$, the following holds for every
  operator $O$ in the image of
  $\cT_{1:i-1}$ applied to the diary operator space:
  \begin{equation}
    \mathbb E\!\left[
      \lVert \Pi_i(O)\rVert_2^2 \,\middle|\, \mathcal F_{i-1}
    \right]
    \le
    (1-\eta)\lVert O\rVert_2^2 .
  \end{equation}
  Then
  \begin{equation}
    \mathbb E\,\lVert \cT_{1:n}(O)\rVert_2^2
    \le
    (1-\eta)^n\lVert O\rVert_2^2
    \le
    e^{-\eta n}\lVert O\rVert_2^2 .
  \end{equation}
  In particular the sequence has an average squared-norm observability
  gap at least $\eta$ on the diary operator space.  This conditional
  contraction is the form expected from physical scrambling dynamics,
  possibly over coarse-grained time windows rather than single record
  steps.

  \item If each one-step algebra is generated by an independent
  uniformly random nonidentity Pauli string on a $D$-dimensional block,
  then for every fixed normalized traceless operator $O$,
  \begin{equation}
    \mathbb E\,
    \lVert \cT_{1:n}(O)\rVert_2^2
    =
    p_q^n,
    \qquad
    p_q=\frac{D^2/2-1}{D^2-1}<\frac12 .
  \end{equation}
  Thus the squared-norm gap is at least $\log 2+O(D^{-2})$, and the
  norm gap is at least $\frac12\log 2+O(D^{-2})$, in expectation.

  \item If the one-step algebras share a common non-scalar commutant
  element $Q$, then every traceless operator in that common commutant has
  $\cT_{1:n}(Q)=Q$ for all $n$.  The observability gap on any diary
  algebra containing such an operator is zero.
\end{enumerate}
\end{proposition}

\begin{proof}
For the first item, iterate the conditional contraction inequality:
\begin{equation}
  \mathbb E\lVert \cT_{1:n}(O)\rVert_2^2
  =
  \mathbb E\,
  \mathbb E\!\left[
    \lVert \Pi_n\cT_{1:n-1}(O)\rVert_2^2
    \middle| \mathcal F_{n-1}
  \right]
  \le
  (1-\eta)\,
  \mathbb E\lVert \cT_{1:n-1}(O)\rVert_2^2 ,
\end{equation}
and repeat.  The second item is the Pauli counting used in
Theorem~\ref{thm:pauli-growth}: a fixed nonidentity Pauli component
survives one random one-step commutant projection with probability
$p_q$, and orthogonality of the Pauli basis gives the claim for any
fixed traceless $O$.  The third item is immediate because a common
commutant operator is fixed by every $\Pi_i$.
\end{proof}

\begin{remark}
The first item is the useful bridge to physical dynamics.  It does not
assume iid records; it asks for a per-step contraction after conditioning
on the past.  Proving such an $\eta$ for Heisenberg iterates
$U^{-i}KU^i$ is the deterministic scrambling problem.  Integrable,
fragmented, symmetry-constrained, or collective-access systems are
expected to fail this condition by leaving $\eta=0$ or parametrically
small.
\end{remark}

\begin{definition}[Effective irreducibility rate]
For records accumulated up to depth $n$, let $\cA_n$ be the algebra
generated by their Heisenberg pullbacks and let $\mathbb E_n$ be the
trace-preserving conditional expectation onto $\cA_n'$.  A block has
effective irreducibility rate $\lambda$ if, for every traceless private
operator $O$ in that block with $\lVert O\rVert_2=1$,
\begin{equation}
  \lVert \mathbb E_n(O)\rVert_2
  \le
  C e^{-\lambda n}.
\end{equation}
\end{definition}

\begin{corollary}[Protected-component collapse]
\label{cor:rate}
If a public block has effective irreducibility rate $\lambda>0$, then
the exact protected component of every normalized traceless private
operator is at most $\varepsilon$ after
\begin{equation}
  n
  \ge
  \lambda^{-1}\log(C/\varepsilon)
\end{equation}
records.
\end{corollary}

\begin{proof}
The conditional expectation $\mathbb E_n$ extracts the component lying
in the commutant $\cA_n'$, which is the exactly protected algebra by
Theorem~\ref{thm:commutant}.  The claim is the definition of effective
irreducibility solved for the record depth $n$.
\end{proof}

This definition isolates the quantitative target.  The structural
commutant theorem says which private information is protected exactly;
an ETH, design, or spectral-gap hypothesis should determine $\lambda$
for a concrete scrambling model.  There are two distinct targets.  For
Hayden-Preskill-type recovery, the relevant quantity is fixed-deposit
de-protection: the protected component of a particular deposited diary
must decay in $O(k+\log(1/\varepsilon))$ independent record units.  The
stronger condition that the entire public block have no protected
subsystem left is worst-case commutant collapse; for a block of
dimension $D$ it can require $O(\log D)$ independent generators, which is
linear in the entropy when $D\sim e^S$.  Collective monitoring has
$\lambda=0$ for both notions.

\begin{theorem}[Random Pauli record growth]
\label{thm:pauli-growth}
This benchmark is the record-algebra analogue of the random Pauli
mechanism behind private quantum channels
\citep{AmbainisMoscaTappDeWolf2000}: independent noncommuting record
generators rapidly erase a fixed protected Pauli component, while exact
worst-case protection persists until the sampled algebra's commutant
collapses.

Let a public block be $q$ qubits, $D=2^q$, and let
$P_1,\ldots,P_n$ be independent uniformly random nonidentity Pauli
strings, modulo phase.  Let $\cA_n$ be the algebra they generate and
let $\mathbb E_n$ be the Hilbert-Schmidt conditional expectation onto
$\cA_n'$.  For every fixed traceless operator $O$ with
$\lVert O\rVert_2=1$,
\begin{equation}
  \mathbb E\,
  \lVert \mathbb E_n(O)\rVert_2^2
  =
  p_q^n,
  \qquad
  p_q =
  \frac{D^2/2-1}{D^2-1}
  < \frac12 .
\end{equation}
Moreover,
\begin{equation}
  \Pr[\cA_n'\ne \mathbb C I]
  \le
  (D^2-1)p_q^n
  \le
  D^2 2^{-n}.
\end{equation}
Thus the fixed-operator protected weight decays exponentially with
record depth, while the worst-case exact protected subsystem is absent
with high probability once
\begin{equation}
  n \ge 2q+\log_2(1/\delta)+O(1).
\end{equation}
\end{theorem}

\begin{proof}
Use the normalized Pauli strings as an orthonormal Hilbert-Schmidt
basis.  The commutant $\cA_n'$ is spanned by the Pauli strings that
commute with every selected $P_i$, so $\mathbb E_n$ keeps exactly those
Pauli components and removes the rest.  For a fixed nonidentity Pauli
string $Q$, the number of nonidentity Pauli strings commuting with $Q$
is $D^2/2-1$, out of $D^2-1$ choices.  Therefore $Q$ survives one random
record generator with probability $p_q$, and survives $n$ independent
generators with probability $p_q^n$.  Expanding a fixed normalized
traceless $O$ in the Pauli basis gives the first claim by linearity of
expectation and orthogonality.

For the second claim, $\cA_n'$ is nontrivial only if some nonidentity
Pauli string commutes with every $P_i$.  A union bound over the
$D^2-1$ nonidentity Pauli strings gives
$(D^2-1)p_q^n$, and $p_q<1/2$ gives the displayed simpler bound.
\end{proof}

\begin{remark}
The theorem separates two notions that are easy to conflate.  A fixed
private operator loses its protected component at an order-one rate per
independent random record generator.  Worst-case protected information,
however, persists as long as the sampled record algebra has a nontrivial
commutant; in this model the exact commutant becomes scalar only after
order $\log D$ independent generators.  This is de-protection, not yet
decodability.  A recovery theorem still needs a channel-capacity,
decoupling, or Hayden-Preskill-type condition.
\end{remark}

\begin{remark}
For a fixed $k$-qubit diary, applying the first estimate to a basis of
$4^k-1$ logical Pauli operators and using Markov plus a union bound gives
fixed-diary de-protection after $O(k+\log(1/\varepsilon)+\log(1/\delta))$
independent record generators, with failure probability $\delta$.  This
is the horizon-relevant scaling.  The iid-Pauli model should be read as a
best-case baseline: in a physical scrambling model the record generators
are typically correlated Heisenberg iterates, such as
$P_i=U^{-i}KU^i$, and the dynamical question is whether chaotic or
ETH-like evolution manufactures effectively independent generators.  The
usual $\log S$ scrambling scale is a time cost for producing such
independent record units, not the number of record units needed once they
are independent.
\end{remark}

\begin{theorem}[Random-coding recovery baseline]
\label{thm:random-recovery}
Let a $k$-qubit diary $D$ be maximally entangled with a reference
$R_D$, and let an old black-hole block $B$ of dimension $d_B$ be
maximally entangled with early radiation $E$.  Apply a Haar-random
unitary, or an exact unitary two-design, to $DB$ and split the output as
new radiation $R$ and remaining block $C$:
\begin{equation}
  U: D B \longrightarrow R C,
  \qquad
  d_D d_B = d_R d_C .
\end{equation}
Then the complementary channel to $C$ is decoupled from the diary
reference on average:
\begin{equation}
  \mathbb E_U
  \left\|
  \rho_{R_D C}
  -
  \rho_{R_D}\otimes\rho_C
  \right\|_1
  \le
  O(d_D/d_R).
\end{equation}
Consequently, if $R$ contains
\begin{equation}
  \log_2 d_R
  \ge
  k+2\log_2(1/\varepsilon)+\log_2(1/\delta)+O(1)
\end{equation}
qubits, then with probability at least $1-\delta$ there is a decoder
$\cR:RE\to \widehat D$ whose entanglement fidelity for the diary is at
least $1-O(\varepsilon)$.
\end{theorem}

\begin{proof}
The trace-norm estimate is the standard decoupling bound used in the
Hayden-Preskill protocol \citep{Hayden2007BlackHoles,Dupuis2014OneShot}:
for a random encoding,
\begin{equation}
  \mathbb E_U
  \left\|
  \rho_{R_D C}
  -
  \rho_{R_D}\otimes\rho_C
  \right\|_1
  \le
  O\!\left(
  \sqrt{\frac{d_D d_C}{d_R d_B}}
  \right)
  =
  O(d_D/d_R),
\end{equation}
where $d_Dd_B=d_Rd_C$ in the last step.  Markov's inequality turns the
average estimate into the stated high-probability bound after increasing
$d_R$ by the displayed overhead.  By the information-disturbance
duality, if the complementary output $C$ is nearly independent of the
diary reference, then the complementary system $RE$ admits a recovery map
for the diary, with entanglement-fidelity loss controlled by the
decoupling error.
\end{proof}

\begin{remark}
This theorem is the access-structure baseline for the fast-recovery
branch: before the record-size threshold the fresh radiation is
unauthorized for the diary, while after the decoupling threshold the
late radiation together with the admissible side information becomes
authorized.  This is standard decoupling physics; the horizon-interface
question is what physical dynamics and coupling geometry manufacture the
threshold at low latency.  Thus the theorem closes the toy baseline:
\begin{equation}
  \text{independent randomization}
  \Longrightarrow
  \text{fixed-diary de-protection}
  \Longrightarrow
  \text{fixed-diary recovery}.
\end{equation}
It also marks the remaining physical problem.  The random-coding
baseline assumes global Haar or design randomization.  In deterministic
source-local models, observability and export are separate tasks: one
must explain how the dynamics produces sufficiently independent
Heisenberg pullbacks of the record coupling, and separately prove the
decoupling condition that makes the emitted records reconstructive.  The
time cost of producing those record units is the scrambling part of the
problem.
\end{remark}

\section{Access Geometry and Export}

The previous two sections isolate two distinct bottlenecks.  The
latency theorem constrains access: before private information reaches
the allowed record channel, the record channel is nearly constant on the
diary.  The decoupling theorem constrains export: once the private
information is in an export-active block, the emitted records recover it
only if the diary reference decouples from the remaining source.
In staged source-local models this gives the recovery window
\begin{equation}
  t_{\rm access}
  \lesssim
  t_{\rm rec}
  \lesssim
  t_{\rm access}+t_{\rm export},
\end{equation}
with the upper end improved when routing, mixing, and export overlap.
The warning is that visibility is not recovery: large commutators,
OTOCs, or boundary sensitivity do not by themselves imply that the
emitted records decouple the diary reference from the remaining source.

\begin{lemma}[Export capacity bound]
\label{lem:export-capacity}
Let $D$ be a $k$-qubit diary maximally entangled with a reference
$R_D$, and let $E$ be admissible side information independent of the
diary at deposit time.  Suppose the only system newly delivered to the
decoder is a record $R$ of dimension $d_R$.  If a decoder on $RE$
recovers the diary with trace-distance error at most $\eta$ on
$R_D\widehat D$, then
\begin{equation}
  \log_2 d_R
  \ge
  k-O\!\left(\eta k+h_2(\eta)\right).
\end{equation}
Consequently, if the record channel exports at most $c_R$ qubits per
step, then
\begin{equation}
  \Delta_{\rm export}
  \ge
  \frac{k-O\!\left(\eta k+h_2(\eta)\right)}{c_R}.
\end{equation}
\end{lemma}

\begin{proof}
Before the record is delivered, admissibility gives $I(R_D:E)=0$.
Sending a system $R$ can increase mutual information with $R_D$ by at
most $2\log_2 d_R$, so
\begin{equation}
  I(R_D:ER)\le 2\log_2 d_R .
\end{equation}
Successful recovery means that the decoded state on $R_D\widehat D$ is
$\eta$-close to a maximally entangled $k$-qubit state.  By continuity of
mutual information,
\begin{equation}
  I(R_D:ER)\ge
  2k-O\!\left(\eta k+h_2(\eta)\right),
\end{equation}
since data processing cannot increase the mutual information under the
decoder.  Combining the two inequalities gives the dimension bound.  The
time bound follows by dividing the required coherent record size by the
per-step export capacity.
\end{proof}

\begin{theorem}[Expander mixer realizes fast-routed recovery]
\label{thm:expander-mixer}
Let $G_S=(V_S,E_S)$ be a bounded-degree graph family with
$|V_S|=S$, and let $X_S\subset V_S$ be the access/emission set.  Suppose
the access radius
\begin{equation}
  R_X(S)=\max_{y\in V_S} d_G(y,X_S)
\end{equation}
is $O(\log S)$, and suppose the source dynamics is finite-speed and
graph-local on $G_S$.

Let a $k$-qubit diary $D$ be deposited in an arbitrary region $Y$, with
reference $R_D$ and admissible side information $E$ independent of the
diary at deposit time.  Assume that for every such $Y$ the diary-bearing
subsystem is collected into an export-active block by time
\begin{equation}
  t_{\rm arr}(Y\to X_S)
  \le
  C_{\rm arr}\log S+O(k).
\end{equation}
Assume also that the subsequent export window emits a record subsystem
$R$ and leaves a remaining source $C$ satisfying
\begin{equation}
  \left\|
  \rho_{R_D C}
  -
  \rho_{R_D}\otimes\rho_C
  \right\|_1
  \le
  \varepsilon .
\end{equation}
Then there is a decoder $\cR:RE\to\widehat D$ recovering the diary with
entanglement-fidelity loss $O(\varepsilon)$, and
\begin{equation}
  t_{\rm rec}
  \le
  C_{\rm arr}\log S+\Delta_{\rm export}.
\end{equation}
If the export window is implemented by a Haar random unitary, an exact
or approximate unitary two-design, a quantum tensor-product expander, or
a random-circuit ensemble satisfying a decoupling theorem, then the
export condition holds with the corresponding theorem's error.  In
particular, once the emitted record has
\begin{equation}
  \log_2 d_R
  \ge
  k+2\log_2(1/\varepsilon)+\log_2(1/\delta)+O(1),
\end{equation}
the standard Hayden-Preskill decoupling budget gives recovery with
failure probability at most $\delta$, up to the error of the chosen
mixer.  Lemma~\ref{lem:export-capacity} gives the corresponding
record-size lower bound: logarithmic access geometry cannot remove the
need to export a diary-sized coherent record, up to continuity terms.
\end{theorem}

\begin{proof}
Before the graph light cone from $Y$ reaches $X_S$, the reduced-record
Lieb-Robinson bound in Lemma~\ref{lem:reduced-lr} makes the diary-to-record
channel diamond-close to a constant channel.  Therefore no decoder on
the records can recover the diary in that interval.

For the upper bound, the assumed collection time places the
diary-bearing subsystem in the export-active block by
$C_{\rm arr}\log S+O(k)$.  The export condition is exactly the
complementary-channel decoupling condition.  By the
information-disturbance duality used in Theorem~\ref{thm:random-recovery},
decoupling of $R_D$ from $C$ implies a recovery map from $RE$ to the
diary.  The displayed record-size threshold is the same
Hayden-Preskill budget as in Theorem~\ref{thm:random-recovery}.  Adding
the collection time and export-window duration gives the stated bound.
\end{proof}

\begin{remark}
There are two readings of the upper bound.  In the clean existence
case, the export-active block is the whole covered source during the
mixer window.  Then there is no separate collection hypothesis: the
diary lies in the mixed block wherever it was deposited, and the
chosen design, tensor-product-expander, or random-circuit primitive
supplies the decoupling condition for a sufficiently large emitted
record.  In the localized-export
variant, where only a bounded region near $X_S$ can emit, collection into
that region is an additional routing hypothesis.  The light cone reaching
$X_S$ gives possible influence on the records; it does not by itself
prove that the diary-bearing subsystem has been collected into the
export-active block.
\end{remark}

\begin{remark}
The theorem is an existence mechanism, not a deterministic Hamiltonian
theorem.  Approximate two-designs and quantum tensor-product expanders
provide theorem-backed export primitives
\citep{Harrow2009RandomCircuits,Harrow2009TPE,Szehr2013Decoupling}.
Brown and Fawzi prove a direct random-circuit decoupling theorem with
$O(n\log^2 n)$ two-qubit gates and $O(\log^3 n)$ depth when the needed
parallel interactions are available \citep{Brown2015Decoupling}.  These
routes establish logarithmic or polylogarithmic fast-routed recovery,
depending on the implementation depth of the chosen mixer.  Replacing
the theorem-backed mixer by a simple fixed nonintegrable Hamiltonian on
an expander graph is the physical strengthening.  Sparse-graph
scrambling results motivate that strengthening, but OTOC growth alone is
not a decoupling/export theorem \citep{Barbon2012Expanders,
Bentsen2019Sparse}.
\end{remark}

\begin{theorem}[Small-subsystem moment-gap export criterion]
\label{thm:moment-gap-export}
Consider the whole-source mixer version of
Theorem~\ref{thm:expander-mixer}.  After the diary is in the mixed
source, let one mixing step be drawn from an ensemble $\nu$ of unitaries
on the active source block, and let
\begin{equation}
  \mathsf M_2(X)
  =
  \mathbb E_{U\sim\nu}
  \left[
    U^{\otimes 2}X(U^\dagger)^{\otimes 2}
  \right]
\end{equation}
be its second-moment channel.  Let $\mathsf P_{\rm Haar}$ denote the
Haar second-moment projector.  Fix the $k$-qubit diary, the emitted
record partition, and the one-shot decoupling functional appearing in
Theorem~\ref{thm:random-recovery}.  Suppose that, in a seminorm
$\|\cdot\|_{\rm dec(D,R)}$ controlling this particular decoupling
functional,
\begin{equation}
  \left\|
    \mathsf M_2^n-\mathsf P_{\rm Haar}
  \right\|_{\rm dec(D,R)}
  \le
  C_{\rm dec}(D,R)e^{-\gamma_2 n}.
\end{equation}
Then, after $n$ mixing steps and emission of a record $R$, the expected
decoupling error obeys
\begin{equation}
  \mathbb E
  \left\|
  \rho_{R_D C}
  -
  \rho_{R_D}\otimes\rho_C
  \right\|_1
  \le
  O(d_D/d_R)+C'_{\rm dec}(D,R)e^{-\gamma_2 n}.
\end{equation}
Consequently, a record satisfying the Hayden-Preskill size budget
together with
\begin{equation}
  n
  \ge
  \gamma_2^{-1}
  \log(C'_{\rm dec}(D,R)/\varepsilon)
\end{equation}
recovers the diary with entanglement-fidelity loss
$O(\varepsilon)$, except with the usual Markov tail probability.  On a
bounded-degree expander with access radius $O(\log S)$, record bandwidth
$c_R$, and $C'_{\rm dec}(D,R)=\exp(O(k+\log(1/\delta)))$, this gives
\begin{equation}
  t_{\rm rec}
  \le
  O(\log S)
  +
  O\!\left(
    \gamma_2^{-1}(k+\log(1/\varepsilon)+\log(1/\delta))
  \right)
  +
  O(k/c_R),
\end{equation}
up to the overheads of the chosen mixing schedule.
\end{theorem}

\begin{proof}
The Hayden-Preskill decoupling estimate in
Theorem~\ref{thm:random-recovery} is a second-moment statement: the Haar
calculation only uses the second moment of the encoding unitary tested
against the diary and the chosen output partition.  If the $n$-step
second-moment channel is within
$C_{\rm dec}e^{-\gamma_2 n}$ of the Haar second moment in a seminorm
that dominates that restricted functional, the same estimate differs
from the Haar bound by at most a constant multiple of this moment error.
This gives
the displayed trace-norm decoupling estimate.  The recovery statement
then follows from the same information-disturbance step used in
Theorem~\ref{thm:random-recovery}.  The final latency bound adds the
expander access time, the moment-mixing time, and the record-size floor
from Lemma~\ref{lem:export-capacity}.
\end{proof}

\begin{remark}
The criterion turns the deterministic expander problem into a precise
gap problem: prove a useful lower bound on $\gamma_2$ for a fixed
bounded-degree expander Hamiltonian or Floquet circuit with the intended
record partition.  The norm in the theorem is deliberately a
small-subsystem/export norm, not a global-design norm.  Global
approximate two-design convergence has an extensive initial deviation
and can require extensive depth; it is stronger than the
Hayden-Preskill recovery task for a fixed diary.  The fast-branch target
is small-subsystem decoupling after the diary has spread through the
expander, with a prefactor controlled by the diary and record budget.
Existing random-circuit results supply nearby benchmarks.  Brown and
Fawzi prove decoupling for random circuits with $O(n\log^2 n)$ gates and
polylogarithmic depth when broad parallel interactions are available
\citep{Brown2015Decoupling}; Harrow and Mehraban prove
short-random-circuit design bounds for local and long-range
architectures \citep{Harrow2023ShortDesigns}, and Mittal--Hunter-Jones
prove approximate-design convergence for random circuits on arbitrary
connected architectures \citep{Mittal2023ArbitraryDesigns}.  These
general architecture results do not by themselves prove
logarithmic-depth export on bounded-degree expanders.  Sparse-graph
fast-scrambling results are evidence for fast operator growth on
expanders, while
Theorem~\ref{thm:moment-gap-export} identifies the stronger
second-moment/export property needed for recovery.
\end{remark}

\begin{remark}
This also fixes the role of the observability and de-protection
diagnostics above.  A positive irreducibility rate $\lambda$ says that
the protected component of a diary operator is being removed relative to
the sampled record algebra.  It does not say that the diary has been
coherently exported.  For example, an expander circuit with fast
controlled-phase spreading and records restricted to a commuting
$Z$-algebra can rapidly remove the protected component of an off-diagonal
diary operator while producing only a dephasing, entanglement-breaking
record channel for that diary.  The observability diagnostic then sees
de-protection, but the Hayden-Preskill decoupling condition fails.
Thus $\lambda$ is a precursor to recovery, while the
second-moment/export gap $\gamma_2$ is the recovery-relevant invariant.
\end{remark}

\section{Horizon-Interface Separation}

The previous results assemble into a family-level separation.  Consider
a sequence of interfaces indexed by a size parameter $S$, with logical
algebra
\begin{equation}
  \cA_{\rm code}^{(S)}
  =
  \bigoplus_x B(\cH_x^{(S)}).
\end{equation}
Call the interface horizon-like, in the access sense used here, if it
has:
\begin{enumerate}
  \item redundant recovery of the sector label $x$ from disjoint passive
  record fragments;
  \item nontrivial private blocks, $\dim\cH_x^{(S)}>1$ for some sectors;
  \item source-side saturation, $N_{\rm eff}(W)=\Theta(S)$ for the
  physically normalized source coupling family; and
  \item location-uniform fast recovery of arbitrary $k$-qubit private
  deposits from the allowed enlarged record/access algebra in
  sub-power-law depth, for example $\operatorname{poly}(k,\log S)$.
\end{enumerate}
The operative obstruction in the theorem is the fourth property, fast
location-uniform private recovery, tested against finite-velocity
source-locality.  The other properties specify the horizon-interface
profile and rule out weaker systems that have only public objectivity,
only private blocks, or only many participating source operators.

\begin{theorem}[Horizon-interface separation]
\label{thm:horizon-interface}
A horizon-like access interface cannot be realized, for all deposit
locations, by finite-velocity source-local emission from a geometrically
local source whose worst-case diary-to-access distance grows as
$L(S)=S^\alpha$ with $\alpha>0$.  Any such realization must contain at
least one of two mechanisms:
\begin{enumerate}
  \item fast internal routing/scrambling, so arbitrary private deposits
  are brought to the emitting/access region in sub-power-law time; or
  \item nonlocal/dressed access, so the allowed record algebra is not the
  reduced algebra generated by source-local emission in the factorized
  source variables.
\end{enumerate}
\end{theorem}

\begin{proof}
By Theorem~\ref{thm:cut}, the redundantly recoverable public part is a
commutative center.  By Corollary~\ref{cor:public-private}, the
noncommuting degrees inside each block cannot themselves be redundantly
public; they are private to the passive fragments.

Assume now that the recovery of arbitrary private deposits is produced
only by finite-velocity source-local emission.  For a worst-case deposit
at distance $L(S)=S^\alpha$ from the access region,
Theorem~\ref{thm:latency} gives recovery latency
$\Omega(S^\alpha)$ up to logarithmic factors, contradicting the assumed
sub-power-law location-uniform recovery.  Theorem~\ref{thm:saturation-slow}
shows that source-side saturation cannot remove this obstruction by
itself, and Theorem~\ref{thm:anonymous-saturation} shows that anonymity,
source symmetry, and unitary record formation still do not force private
mixing.  Therefore a successful interface must either remove the
power-law routing bottleneck by fast internal mixing, or change the
accessible algebra so that the records are nonlocal/dressed relative to
the factorized source variables.
\end{proof}

\begin{remark}
The theorem classifies mechanisms, not ontologies.  A single physical
system can admit both descriptions: a stretched-horizon account may
emphasize fast routing, while a boundary or gravitationally dressed
account may emphasize the nonlocal access algebra.
\end{remark}

\section{Examples and Stress Tests}

Ordinary local reservoirs fall on the slow side of
Theorem~\ref{thm:latency}.  Theorem~\ref{thm:saturation-slow} shows that
source-side participation $N_{\rm eff}(W)=\Theta(S)$ does not by itself
imply logarithmic latency.  Theorem~\ref{thm:anonymous-saturation}
shows that even anonymous symmetric unitary access can stop at a public
collective center and leave private blocks unmixed; Theorem~\ref{thm:commutant}
identifies those unmixed blocks as the commutant/noiseless factors of
the record-generating algebra.  Theorem~\ref{thm:expander-mixer} gives a
clean positive fast-routing witness: a bounded-degree log-diameter
access geometry, combined with a theorem-backed export mixer, has
location-uniform logarithmic or polylogarithmic private recovery.  In
fast-scrambler and stretched-horizon pictures, the required role of the
dynamics is to provide both rapid access and coherent export, with a
floor set by scrambling plus the emitted-record budget.  Nonlocal
encoders, boundary reconstruction, and gravitationally dressed access
realize the nonlocal/dressed branch; by Corollary~\ref{cor:frozen}, this
branch is diagnosed by recovery that survives after source-local routing
is frozen.  Rindler wedges test thermality and observer
algebras but do not by themselves supply a finite-sector Page/Hayden-
Preskill recovery channel.  BTZ/AdS examples can realize both branches
depending on whether the protocol is described as boundary
reconstruction or as chaotic routing in the boundary/stretched-horizon
dynamics.

\section{Relation to Horizon Interfaces}

The cut theorem supplies the public/classical side of a proposed
constrained-access horizon interface: redundant recoverability selects
the center.  The latency theorem supplies the private/quantum side:
block information cannot be recovered quickly from source-local records
unless routing is fast or the access algebra is nonlocal/dressed.
The expander mixer theorem supplies the positive fast-routing side:
logarithmic access geometry plus decoupling export realizes
location-uniform fast private recovery.  Theorem~\ref{thm:horizon-interface}
combines these into the access profile:
\begin{equation}
\begin{gathered}
  \text{public center from redundant records,}\\
  \text{private blocks hidden from small/passive records,}\\
  \text{fast private recovery by log-depth access plus export,}\\
  \text{or by dressed/nonlocal access.}
\end{gathered}
\end{equation}
This separates horizon-like private recovery from ordinary thermality,
ordinary Darwinian records, and direct high-bandwidth monitoring.

\bibliographystyle{unsrtnat}
\bibliography{refs}

\end{document}
